\title{DeepSeekMath: Pushing the Limits of Mathematical Reasoning in Open Language Models}

\begin{abstract}
The intricacies and structured demands of mathematical reasoning present notable difficulties for language models. In this study, we present DeepSeekMath~7B, an extension of DeepSeek-Coder-Base-v1.5 7B through continued pre-training using 120B tokens associated with mathematics, drawn primarily from Common Crawl, alongside corpora of natural language and code. DeepSeekMath~7B attains a notable 51.7\% accuracy on the MATH benchmark at the competition level, all without the support of auxiliary toolkits or ensemble approaches, and achieving results close to those of Gemini-Ultra and GPT-4. When applying self-consistency across 64 generated samples, DeepSeekMath~7B attains 60.9\% accuracy on MATH. The advancement of DeepSeekMath~in mathematical reasoning is attributable to two main elements: (1) the exploitation of extensive, publicly-available web sources through a carefully designed data selection methodology, and (2) the development of Group Relative Policy Optimization (GRPO)—an adaptation of Proximal Policy Optimization (PPO)—which augments the model’s reasoning abilities in mathematics and yields improvements in PPO memory efficiency.
\end{abstract}

\section{Introduction}

Large language models (LLMs) have transformed how artificial intelligence tackles mathematical reasoning, leading to notable progress on benchmarks for quantitative reasoning \citep{MATH} as well as for geometry \citep{trinh2024solving}. In addition, these systems have shown great utility in supporting humans as they work through challenging mathematical questions \citep{tao}. Nevertheless, state-of-the-art systems like GPT-4 \citep{gpt4} and Gemini-Ultra \citep{gemini} remain inaccessible to the public, while open-source alternatives currently lag substantially in terms of performance.

In this work, we present DeepSeekMath, a specialized language model focused on mathematics that considerably surpasses existing open-source alternatives in mathematical reasoning and closely rivals GPT-4 on standard academic evaluations. Our approach centers on constructing the DeepSeekMath Corpus, an extensive pre-training dataset containing 120B math tokens, meticulously curated for quality. This collection is sourced from the Common Crawl (CC) using a classifier based on fastText \citep{joulin2016fasttext}. Initially, the classifier is trained with positive samples from OpenWebMath \citep{openwebmath} and negative samples drawn from a broad array of unrelated web content. Following this, we use the classifier to extract new positive math-related data from the CC, which are subsequently validated and enhanced by human annotators. The classifier is then refined using this augmented dataset, resulting in improved identification accuracy. Evaluations demonstrate the high caliber of our large-scale corpus: DeepSeekMath-Base 7B attains 64.2\% accuracy on GSM8K \citep{gsm8k} and 36.2\% on the MATH competition dataset \citep{MATH}, exceeding the results of Minerva 540B \citep{minerva}. Furthermore, since the DeepSeekMath Corpus is multilingual, we also observe progress on Chinese mathematical benchmarks \citep{wei2023cmath,agieval}. We regard our methodology in mathematical data collection and preparation as a foundation for ongoing advancements in this field, with substantial opportunities for further improvements ahead.

DeepSeekMath-Base is built upon DeepSeek-Coder-Base-v1.5 7B \citep{deepseek-coder}, as we find that initializing with a code-focused model leads to superior results relative to beginning with a standard LLM. Additionally, our findings show that mathematical instruction not only boosts the model’s performance on tasks requiring mathematics, but also strengthens outcomes on broader evaluations such as MMLU \citep{mmlu} and BBH \citep{bbh}. This suggests that math-oriented training contributes to enhanced general reasoning skills beyond its intended domain.

Following the pre-training phase, we fine-tune DeepSeekMath-Base on mathematical instructions utilizing datasets featuring chain-of-thought \citep{cot}, program-of-thought \citep{pot,pal}, and tool-integrated reasoning \citep{tora} approaches. This procedure yields the DeepSeekMath-Instruct 7B model, which surpasses all other 7B models and achieves performance on par with 70B-scale open-source models that have been instruction-tuned.

Additionally, we present Group Relative Policy Optimization (GRPO), a novel reinforcement learning (RL) approach inspired by Proximal Policy Optimization (PPO) \citep{schulman2017proximal}. Unlike traditional PPO, GRPO eliminates the need for a critic network, instead deriving the baseline directly from group-level scores, thereby substantially lowering the computational demands during training. Utilizing only a fraction of English instruction tuning data, GRPO yields marked improvements over the already competitive DeepSeekMath-Instruct, both on in-domain tasks (GSM8K: 82.9\% $\rightarrow$ 88.2\%, MATH: 46.8\% $\rightarrow$ 51.7\%) as well as on out-of-domain mathematical benchmarks (for instance, CMATH: 84.6\% $\rightarrow$ 88.8\%) in the RL training stage. Furthermore, we develop a comprehensive framework to analyze and relate various approaches, including Rejection Sampling Fine-Tuning (RFT) \citep{yuan2023scaling}, Direct Preference Optimization (DPO) \citep{dpo}, PPO, and GRPO. Through this lens, it becomes clear that these methodologies can be categorized broadly as direct or simplified RL strategies. We perform a series of rigorous experiments—such as comparisons between online and offline learning, distinctions between outcome-based and process-based supervision, and analyses of single-turn versus iterative RL—to explore the core facets of this unified paradigm in depth. Finally, we elucidate the mechanisms by which our RL methodology enhances instruction-tuned models’ performance and propose prospective avenues for advancing RL techniques within this conceptual framework.

\subsection{Contributions}
We present scalable approaches to math pre-training and investigate reinforcement learning methods through detailed analysis and experimentation.

\textbf{Math Pre-Training at Scale}

\begin{itemize}[topsep=0pt]
    \item Our study demonstrates that the widely available Common Crawl dataset includes significant mathematical content. Through the use of a carefully engineered data selection framework, we have assembled the DeepSeekMath Corpus, comprising 120B tokens from web sources rigorously filtered for mathematical relevance. This dataset is nearly seven times larger than the math-specific web data employed by Minerva \citep{minerva}, and surpasses the recently released OpenWebMath \citep{openwebmath} by a factor of nine.
\end{itemize}

\item DeepSeekMath-Base 7B, our pre-trained foundational model, attains results similar to those of Minerva 540B \citep{minerva}, suggesting that parameter count alone does not solely determine proficiency in mathematical reasoning. Robust outcomes can also be realized by a smaller model trained on superior-quality datasets.

\item We present results from our experiments on mathematical training. Engaging in code training before initiating math-specific training enhances the model’s proficiency in addressing mathematical tasks, regardless of whether tool use is involved. These results contribute to the ongoing debate regarding the impact of code training on reasoning capabilities, suggesting that such training does enhance reasoning—specifically in mathematical contexts.

\item While it is typical to use arXiv papers for training, particularly within mathematics-focused studies, this approach does not yield significant gains across any of the mathematical benchmarks considered in this work.
\end{itemize}

\textbf{Exploration and Analysis of Reinforcement Learning}

\begin{itemize}[topsep=0pt]
    \item We present Group Relative Policy Optimization (GRPO), a novel reinforcement learning approach that operates without the need for a critic model. By deriving the baseline from group-based scores, GRPO achieves substantial reductions in computational and training demands compared to Proximal Policy Optimization (PPO).
    \item Our experiments reveal that applying GRPO to DeepSeekMath-Instruct—exclusively utilizing instruction-tuning data—yields notable improvements in its capabilities. Additionally, we find that GRPO boosts out-of-domain generalization throughout the reinforcement learning phase.
    \item We propose a comprehensive framework that contextualizes methods such as RFT, DPO, PPO, and GRPO under a common lens. A suite of empirical studies is carried out, including comparisons between online and offline regimes, supervision via outcomes versus process, and analyses of single-turn against iterative reinforcement learning, to rigorously examine the critical aspects underpinning this framework.
    \item Leveraging our integrative framework, we analyze the mechanisms that underpin the success of reinforcement learning, and outline several promising avenues for enhancing the effectiveness of reinforcement learning in large language models (LLMs).
\end{itemize}

\subsection{Summary of Evaluations and Metrics}

\begin{itemize}[topsep=0pt]
    \item \textbf{English and Chinese Mathematical Reasoning}: Our evaluation spans both English and Chinese benchmarks, addressing mathematical challenges from elementary through collegiate levels. The English datasets comprise GSM8K \citep{gsm8k}, MATH \citep{MATH}, SAT \citep{llemma}, OCW Courses \citep{minerva}, and MMLU-STEM \citep{mmlu}. For Chinese, we utilize MGSM-zh \citep{mgsm}, CMATH \citep{wei2023cmath}, Gaokao-MathCloze \citep{agieval}, and Gaokao-MathQA \citep{agieval}. We assess models both on their capability to produce fully self-contained written solutions independently of external tools, and on their proficiency in using Python to arrive at correct answers.

When evaluated on English benchmarks, DeepSeekMath-Base performs on par with the proprietary Minerva 540B model \citep{minerva} and outperforms all other available open-source base models, including Mistral 7B \citep{mistral} and Llemma-34B \citep{llemma}, regardless of their exposure to math-specific pre-training, frequently by a considerable margin. Importantly, DeepSeekMath-Base also demonstrates superior results on Chinese benchmarks, which may be attributed to our approach of not restricting math pre-training data to English, as done in prior studies \citep{minerva,llemma}, and instead incorporating high-quality data in various non-English languages. With the implementation of mathematical instruction tuning and reinforcement learning, the advanced models DeepSeekMath-Instruct and DeepSeekMath-RL exhibit impressive capabilities, achieving over 50\% accuracy on the competition-level MATH dataset—representing a first among open-source models.

\item \textbf{Formal Mathematics}: We assess the performance of DeepSeekMath-Base on the task of informal-to-formal theorem proving as outlined in \citep{dsp_proof}, utilizing miniF2F \citep{minif2f} and selecting Isabelle \citep{isabelle} as the proof assistant. Our findings indicate that DeepSeekMath-Base achieves notable results in few-shot autoformalization scenarios.
\item \textbf{Natural Language Understanding, Reasoning, and Code}: To establish an in-depth evaluation of the models’ abilities in language comprehension, logical reasoning, and code synthesis, DeepSeekMath-Base is benchmarked against the Massive Multitask Language Understanding (MMLU) dataset \citep{mmlu}, which consists of 57 diverse multiple-choice tests, BIG-Bench Hard (BBH) \citep{bbh}, a collection of 23 complex reasoning tasks, and coding assessments including HumanEval \citep{codex} and MBPP \citep{mbpp}, both of which are standard for gauging programming language model proficiency. Pre-training on mathematical content proves advantageous for both linguistic understanding and reasoning proficiency.

\section{Math Pre-Training}

\subsection{Data Collection and Decontamination} Here, we describe the methodology used to assemble the DeepSeekMath Corpus utilizing data from Common Crawl. As illustrated in Figure \ref{fig:data_collect}, an iterative framework is employed, systematically harvesting a comprehensive mathematical dataset beginning with an initial seed collection composed of a limited but curated set of math-oriented data. It should be emphasized that this procedure is not limited to mathematics; it can be readily adapted to domains like programming as well.

Initially, we select OpenWebMath \citep{openwebmath}, an assemblage of well-curated mathematical web documents, to serve as our seed dataset. With this foundation, we train a fastText model \citep{joulin2016fasttext} designed to identify additional mathematical webpages resembling those in OpenWebMath. Our methodology involves sampling 500,000 items from the seed corpus as positive examples, while an equal number of webpages from Common Crawl act as negative samples. Training is conducted using a publicly available library\footnote{\url{https://fasttext.cc}}, employing a vector size of 256, a learning rate of 0.1, maximum word n-gram length set to 3, a minimum of 3 word occurrences, and 3 training epochs. To address the substantial volume of Common Crawl data, we apply deduplication both by URL and using near-duplicate detection, ultimately producing a corpus of 40 billion HTML pages. The trained fastText model is subsequently used to retrieve mathematical web pages from the deduplicated Common Crawl data. To eliminate inferior-quality content, the resulting pages are ranked by the confidence scores produced by the model, with only the top-ranking documents retained. We determine the appropriate data scale by performing pre-training trials on the highest-scoring subsets comprising 40B, 80B, 120B, and 160B tokens. For the initial round, we opt to retain the top 40B tokens.

Following the initial round of data collection, a significant portion of mathematical web pages remains uncollected. This shortfall primarily results from the fastText model being trained on a set of positive samples that lacks adequate variety. To address this, we broaden the seed corpus by sourcing additional mathematical web resources, thereby allowing us to refine the fastText model further. Our strategy begins by dividing the complete Common Crawl dataset into non-overlapping domains, where each domain comprises web pages sharing a common base URL. For each domain, we determine the fraction of web pages gathered during the initial iteration. If over 10\% of a domain’s web pages are collected, we designate it as math-related (for example, \url{mathoverflow.net}).

Next, we perform manual annotation to identify URLs within these math-related domains that are associated with mathematical topics (such as \url{mathoverflow.net/questions}). Uncollected web pages that correspond to these annotated URLs are then incorporated into the seed corpus. By doing so, we are able to obtain a broader set of positive instances, enabling us to train a more effective fastText model with improved recall for mathematical content in the subsequent data collection phase. After four rounds of this iterative process, we accumulate 35.5M mathematical web pages, comprising a total of 120B tokens. During the fourth iteration, it becomes apparent that approximately 98\% of the data had already been retrieved in the third round, which prompts us to discontinue further collection.

In order to prevent benchmark contamination, we adopt the approach described by \cite{deepseek-coder}, which involves filtering out web pages that include questions or answers from English mathematical benchmarks such as GSM8K \citep{gsm8k} and MATH \citep{MATH}, as well as Chinese benchmarks like CMATH \citep{wei2023cmath} and AGIEval \citep{agieval}. The filtering process operates as follows: if a segment of text contains a 10-gram sequence that matches exactly any substring from the evaluation benchmarks, that segment is excluded from our mathematical training dataset. For benchmark passages with fewer than 10 but at least 3 grams, we apply the same exact-match strategy to eliminate any potentially contaminated web pages.

\subsection{Validating the Quality of the DeepSeekMath~Corpus}

We run pre-training experiments to investigate how the DeepSeekMath~Corpus is compared with the recently released math-training corpora:

\begin{itemize}[topsep=0pt]
    \item \textbf{MathPile} \citep{mathpile}: an extensive dataset comprising 8.9B tokens drawn from a variety of sources such as textbooks, Wikipedia, ProofWiki, CommonCrawl, StackExchange, and predominantly arXiv (contributing more than 85\% of the data);
    \item \textbf{OpenWebMath} \citep{openwebmath}: a 13.6B-token dataset obtained by filtering CommonCrawl for mathematical material;
    \item \textbf{Proof-Pile-2} \citep{llemma}: this resource integrates OpenWebMath, AlgebraicStack (which supplies 10.3B tokens of mathematical code), and arXiv publications (adding 28.0B tokens). In experiments with Proof-Pile-2, we maintain the arXiv:Web:Code ratio of 2:4:1, as specified in \cite{llemma}.
\end{itemize}

\subsubsection{Training Setting}
\label{sec:quality-policy}
We conduct mathematical training on a broadly pre-trained language model containing 1.3B parameters, based on the same architecture as the DeepSeek LLMs \citep{deepseek-llm}, referred to as DeepSeek-LLM 1.3B. The model is individually trained on each mathematical dataset for a total of 150B tokens. All training procedures leverage the HAI-LLM framework \citep{haillm}, known for its efficiency and lightweight design. Adhering to the methodology used by DeepSeek LLMs, we utilize the AdamW optimizer \citep{adamW} with $\beta_1=0.9$, $\beta_2=0.95$, and a $\mathrm{weight\_decay}$ set to 0.1. The learning rate is managed by a multi-stage scheduler: it ramps up to its maximum after 2,000 warmup steps, decays to 31.6\% of its maximum by 80\% of the training duration, and subsequently drops to 10.0\% of the peak by the 90\% mark. The highest learning rate applied is 5.3e-4, with training performed using a batch size of 4 million tokens and a context window of 4K.

\subsubsection{Evaluation Results}

\textbf{The DeepSeekMath~Corpus is of high quality, covers multilingual mathematical content, and is the largest in size.}

\begin{itemize}[topsep=0pt]
    \item \textbf{High-quality}:
    To assess downstream capabilities, we measure performance across 8 mathematical benchmarks employing few-shot chain-of-thought prompting \cite{cot}. Table \ref{tab:corpora_comparison} highlights that models trained with the DeepSeekMath~Corpus consistently achieve superior results. Furthermore, Figure \ref{fig:corpora_comparisons} illustrates that the DeepSeekMath Corpus outperforms Proof-Pile-2 at the 50B token mark (equivalent to a single epoch over Proof-Pile-2), providing evidence that the DeepSeekMath Corpus offers higher average data quality.
    \item \textbf{Multilingual}:
    The DeepSeekMath~Corpus contains data spanning several languages, with English and Chinese as its predominant components. Table \ref{tab:corpora_comparison} indicates that models trained on this corpus experience improved mathematical reasoning proficiency in both English and Chinese. By contrast, existing math corpora—mostly focused on English—result in minimal improvement or even adverse effects when evaluated on Chinese mathematical reasoning tasks.
    \item \textbf{Large-scale}:
    The DeepSeekMath~Corpus surpasses prior mathematical corpora by a considerable margin in size. According to Figure \ref{fig:corpora_comparisons}, DeepSeek-LLM 1.3B, when trained using the DeepSeekMath~Corpus, exhibits accelerated and more sustained performance gains. In comparison, baseline corpora are significantly smaller and are cycled through the training process multiple times, resulting in models whose progress soon stagnates.
\end{itemize}

\subsection{Training and Evaluating DeepSeekMath-Base 7B}

In this section, we present DeepSeekMath-Base 7B, a foundational model that demonstrates robust mathematical reasoning skills. Building on DeepSeek-Coder-Base-v1.5 7B \citep{deepseek-coder} as its initialization point, the model undergoes training on 500B tokens. The data is allocated as follows: 56\% sourced from the DeepSeekMath Corpus, 4\% derived from AlgebraicStack, 10\% from arXiv publications, 20\% comprising Github code, and the final 10\% consisting of natural language content taken from Common Crawl in both English and Chinese. The model training primarily follows the procedures outlined in Section \ref{sec:quality-policy}, with two notable differences: the learning rate peaks at 4.2e-4, and a token batch size of 10M is used.

We thoroughly evaluate DeepSeekMath-Base 7B to determine its mathematical proficiency, specifically analyzing its capacity to independently generate complete mathematical solutions, apply tools to resolve mathematical problems, and engage in formal theorem proving. In addition to examining its mathematical functionality, we present a broader characterization of the base model by reporting its capabilities in natural language understanding, reasoning, and programming.

\paragraph{Mathematical Problem Solving with Step-by-Step Reasoning}

We assess the mathematical problem-solving abilities of DeepSeekMath-Base utilizing few-shot chain-of-thought prompting \citep{cot} on eight distinct benchmarks in both English and Chinese. These datasets include a range of quantitative reasoning tasks—such as GSM8K \citep{gsm8k}, MATH \citep{MATH}, and CMATH \citep{wei2023cmath}—as well as multiple-choice assessments like MMLU-STEM \citep{mmlu} and Gaokao-MathQA \citep{agieval}. Collectively, they span various mathematical domains and difficulties, from primary school arithmetic to college-level topics.

According to the data presented in Table \ref{tab:base_math_cot}, DeepSeekMath-Base 7B achieves the highest performance on all eight evaluated benchmarks among open-source base models. This includes outperforming prominent models such as the general-purpose Mistral 7B \citep{mistral} and the recently introduced Llemma 34B \citep{llemma}, which has received additional mathematical training on Proof-Pile-2 \citep{llemma}. Particularly noteworthy is DeepSeekMath-Base’s result on the competition-level MATH dataset, where it exceeds the leading open-source base models by more than 10 percentage points, and also achieves superior results compared to Minerva 540B \citep{minerva}—a closed-source base model based on PaLM \citep{palm}, 77 times larger and trained with mathematical corpora.

\paragraph{Mathematical Problem Solving with Tool Use} We assess the effectiveness of program-assisted mathematical reasoning on GSM8K and MATH datasets through the use of few-shot program-of-thought prompting \citep{pot,pal}. The models are instructed to generate a Python script to address each question, making use of available libraries such as \textit{math} and \textit{sympy} to carry out complex calculations. The program's output is then treated as the solution to the problem. According to Table \ref{tab:base_math_pot_proof}, DeepSeekMath-Base 7B achieves higher performance compared to the previous leading model, Llemma 34B.

\paragraph{Formal Mathematics}

Automating formal proofs has become increasingly prominent due to its role in bolstering the correctness and dependability of mathematical arguments while streamlining the proving process. In this work, we assess DeepSeekMath-Base 7B on the informal-to-formal proof generation task described in \citep{dsp_proof}, wherein the objective is to construct a formal proof from an informal mathematical statement, its formal representation, and a corresponding informal proof. Our evaluation employs the miniF2F benchmark \citep{minif2f}, which consists of challenging Olympiad-style mathematical problems. For each task, we use few-shot prompting to generate a formal Isabelle proof. Adhering to the methodology of \cite{dsp_proof}, we prompt language models to create proof sketches and subsequently invoke Sledgehammer \citep{sledgehammer}, a standard automated prover, to complete any gaps. The results summarized in Table \ref{tab:base_math_pot_proof} indicate that DeepSeekMath-Base 7B achieves competitive results in automating the formalization of proofs.

\paragraph{Natural Language Understanding, Reasoning, and Code}

We assess the abilities of our models in natural language understanding using the MMLU \citep{mmlu} benchmark, evaluate reasoning with BBH \citep{bbh}, and measure coding proficiency on HumanEval \citep{codex} and MBPP \citep{mbpp}. According to the data presented in Table \ref{tab:understanding_reasoning_code}, DeepSeekMath-Base 7B demonstrates notable improvements in both MMLU and BBH tasks compared to the earlier DeepSeek-Coder-Base-v1.5 model \citep{deepseek-coder}, highlighting the benefits of incorporating mathematical pre-training for language understanding and reasoning capabilities. Furthermore, the integration of code tokens during continual training allows DeepSeekMath-Base 7B to retain the coding performance established by DeepSeek-Coder-Base-v1.5 across both coding assessments. Collectively, DeepSeekMath-Base 7B achieves substantial gains over the general-purpose model Mistral 7B \citep{mistral} on all three reasoning and coding benchmarks.

\section{Supervised Fine-Tuning}

\subsection{SFT Data Curation}

We develop a dataset for instruction tuning in mathematics that encompasses both English and Chinese questions, spanning multiple areas of mathematics and a range of difficulty. Each problem is matched with a corresponding solution presented in chain-of-thought (CoT) \citep{cot}, program-of-thought (PoT) \citep{pot,pal}, or tool-integrated reasoning style \citep{tora}. The dataset includes a total of 776K examples for training.

\begin{itemize}[topsep=0pt]
    \item \textbf{English mathematical datasets}: We provide tool-integrated annotations for GSM8K and MATH problem sets. In addition, we include selected examples from MathInstruct \citep{MathInstruct} as well as the training portion of Lila-OOD \citep{lila}, both featuring solutions utilizing CoT or PoT approaches. This English dataset compilation encompasses a wide range of mathematical domains, including but not limited to algebra, probability, number theory, calculus, and geometry.
    \item \textbf{Chinese mathematical datasets}: We assemble Chinese K-12 math problems that span 76 distinct sub-domains, such as linear equations, each annotated with solutions in both CoT and tool-augmented reasoning formats.
\end{itemize}

\subsection{Training and Evaluating DeepSeekMath-Instruct 7B}

This section presents DeepSeekMath-Instruct 7B, a model derived from DeepSeekMath-Base through instruction tuning on mathematical tasks. During training, examples are randomly joined together up to a maximum context window of 4K tokens. The model is optimized over 500 training steps, utilizing a batch size of 256 and maintaining a fixed learning rate of 5e-5 throughout.

We evaluate models' mathematical performance both without and with tool use, on 4 quantitative reasoning benchmarks in English and Chinese. We benchmark our model against the leading models of the time:

\begin{itemize}[topsep=0pt]
    \item \textbf{Closed-source models} encompass: (1) the GPT series, with GPT-4 \citep{gpt4} and GPT-4 Code Interpreter~\footnote{\url{https://openai.com/blog/chatgpt-plugins\#\#code-interpreter}} representing the most advanced examples, (2) Gemini Ultra and Pro \citep{gemini}, (3) Inflection-2 \citep{inflection-2}, (4) Grok-1~\footnote{\url{https://x.ai/model-card}}, along with recent releases by Chinese firms such as (5) Baichuan-3~\footnote{\url{https://www.baichuan-ai.com}}, and (6) the newest iteration, GLM-4~\footnote{\url{https://open.bigmodel.cn/dev/api\#glm-4}}.

    from the GLM family \citep{glm}.

These models are intended for broad applications and have generally been subjected to multiple alignment processes. 
\item \textbf{Open-source models} encompass: general-purpose models such as (1) DeepSeek-LLM-Chat 67B \citep{deepseek-llm}, (2) Qwen 72B \citep{qwen}, (3) SeaLLM-v2 7B \citep{seallm}, and (4) ChatGLM3 6B \citep{chatglm3}. In addition, there are models particularly tailored for mathematical capabilities, including (5) InternLM2-Math 20B~\footnote{\url{https://github.com/InternLM/InternLM-Math}}, which is derived from InternLM2, having received specialized training on mathematical data before undergoing instruction tuning; (6) Math-Shepherd-Mistral 7B, which utilizes PPO training \citep{schulman2017proximal} on Mistral 7B \citep{mistral} coupled with a process-supervised reward model; (7) the WizardMath series \citep{wizardmath}, which enhances the mathematical reasoning ability of Mistral 7B and Llama-2 70B \citep{llama2} through evolve-instruct (an instruction tuning technique leveraging AI-generated instructions) and PPO training, predominantly with questions sourced from GSM8K and MATH; (8) MetaMath 70B \citep{metamath}, which is Llama-2 70B further adapted using a version of GSM8K and MATH that has been extended; (9) ToRA 34B \cite{tora}, constructed by fine-tuning CodeLlama 34B to enable mathematical problem solving with tool integration; and (10) MAmmoTH 70B \citep{MathInstruct}, which involves instruction tuning of Llama-2 70B using MathInstruct.

As indicated in Table \ref{tab:sft_rl_math}, when evaluated without the use of external tools, DeepSeekMath-Instruct 7B excels at step-by-step reasoning. Specifically, within the competition-tier MATH dataset, our model outperforms all existing open-source alternatives, as well as most proprietary systems such as Inflection-2 and Gemini Pro, achieving an improvement of at least 9\% in absolute terms. This advantage persists even when compared to models with considerably larger parameter counts (for example, Qwen 72B) or those that have received additional fine-tuning via mathematics-oriented reinforcement learning strategies (such as WizardMath-v1.1 7B). Although DeepSeekMath-Instruct achieves performance on par with Chinese proprietary models GLM-4 and Baichuan-3 on MATH, its results do not yet reach the level demonstrated by GPT-4 or Gemini Ultra.

When evaluated in a scenario that permits the combination of natural language reasoning with programmatic tool utilization, DeepSeekMath-Instruct 7B attains nearly 60\% accuracy on the MATH benchmark, exceeding the performance of all previously released open-source models. Additionally, on other evaluation datasets, our model achieves results comparable to DeepSeek-LLM-Chat 67B, the former leading approach despite being an order of magnitude larger.

\subsection{Group Relative Policy Optimization} Recent studies have demonstrated that reinforcement learning (RL) can significantly enhance the mathematical reasoning skills of LLMs following the Supervised Fine-Tuning (SFT) phase \citep{wang2023math,wizardmath}. In the following, we present Group Relative Policy Optimization (GRPO), our proposed RL algorithm that is both efficient and effective.

\subsubsection{From PPO to GRPO} Proximal Policy Optimization (PPO) \citep{schulman2017proximal} is an actor-critic RL algorithm that is widely used in the RL fine-tuning stage of LLMs \citep{ouyang2022training}. In particular, it optimizes LLMs by maximizing the following surrogate objective:

\begin{equation}
\footnotesize
    \mathcal{J}_{PPO}(\theta) = \mathbb{E}{[q \sim P(Q),\; o \sim \pi_{\theta_{old}}(O|q)]} \frac{1}{|o|} \sum_{t=1}^{|o|} \min \Bigg\{ \frac{\pi_\theta(o_{t} | q, o_{<t})}{\pi_{\theta_{old}}(o_{t} | q, o_{<t})} A_{t},\; \text{clip}\Big( \frac{\pi_\theta(o_{t} | q, o_{<t})}{\pi_{\theta_{old}}(o_{t} | q, o_{<t})}, 1 - \varepsilon, 1 + \varepsilon \Big)  A_{t} \Bigg\} ,
\end{equation}

Here, $\pi_{\theta}$ denotes the current policy model, while $\pi_{\theta_{old}}$ refers to the policy model prior to the most recent update. The variables $q$ and $o$ represent questions and their corresponding outputs, respectively, which are drawn from the question dataset and produced according to the previous policy $\pi_{\theta_{old}}$. The hyper-parameter $\varepsilon$ is used to control clipping during training, as introduced in PPO, to ensure training stability. The term $A_t$ stands for the advantage estimate, which is calculated utilizing Generalized Advantage Estimation (GAE) \citep{gae}, relying on the future rewards $\{r_{\ge t}\}$ and an estimated value function $V_{\psi}$. Consequently, PPO requires simultaneous optimization of a value function in addition to the policy. To further prevent the reward model from being excessively optimized, it is common practice to apply a KL penalty at every token, which is computed relative to a reference model and incorporated into the reward, as outlined by \citep{ouyang2022training}, that is,

\begin{equation}
    r_{t} = r_\varphi(q, o_{\le t}) - \beta \log\left(\frac{\pi_{\theta}(o_{t}|q, o_{<t})}{\pi_{ref}(o_{t}|q, o_{<t})}\right),
\label{eq:PPO-reward}
\end{equation}

Here, $r_\varphi$ denotes the reward model, $\pi_{ref}$ represents the reference policy—commonly set to the original SFT model—and $\beta$ serves as the weight applied to the KL divergence penalty.

As the value function employed in PPO is typically another model of comparable size as the policy model, it brings a substantial memory and computational burden. Additionally, during RL training, the value function is treated as a baseline in the calculation of the advantage for variance reduction. While in the LLM context, usually only the last token is assigned a reward score by the reward model, which may complicate the training of a value function that is accurate at each token. To address this, as shown in Figure \ref{fig:grpo}, we propose Group Relative Policy Optimization (GRPO), which obviates the need for additional value function approximation as in PPO, and instead uses the average reward of multiple sampled outputs, produced in response to the same question, as the baseline. More specifically, for each question $q$, GRPO samples a group of outputs $\{o_1, o_2, \cdots, o_G\}$  from the old policy  $\pi_{\theta_{old}}$  and then optimizes the policy model by maximizing the following objective:

\begin{equation}
\footnotesize
\begin{split}
    \mathcal{J}_{GRPO}(\theta) &= \mathbb{E}{[q \sim P(Q), \{o_i\}_{i=1}^G \sim \pi_{\theta_{old}}(O|q)]}  \\
    & \frac{1}{G}\sum_{i=1}^G\frac{1}{|o_i|} \sum_{t=1}^{|o_i|} \left\{ \min \left[ \frac{\pi_\theta(o_{i,t} | q, o_{i,<t})}{\pi_{\theta_{old}}(o_{i,t} | q, o_{i,<t})} \hat{A}_{i,t}, \text{clip} \left( \frac{\pi_\theta(o_{i,t} | q, o_{i,<t})}{\pi_{\theta_{old}}(o_{i,t} | q, o_{i,<t})}, 1 - \varepsilon, 1 + \varepsilon \right)  \hat{A}_{i,t} \right] - \beta \mathbb{D}_{KL}\left[\pi_{\theta} || \pi_{ref}\right]\right\} ,
\end{split}
\label{eq:GRPO-obj}
\end{equation}

Rephrased:

\begin{equation}
\footnotesize
\begin{split}
    \mathcal{J}_{GRPO}(\theta) &= \mathbb{E}{[q \sim P(Q), \{o_i\}_{i=1}^G \sim \pi_{\theta_{old}}(O|q)]}  \\
    & \frac{1}{G}\sum_{i=1}^G\frac{1}{|o_i|} \sum_{t=1}^{|o_i|} \bigg\{ \min \bigg[ \frac{\pi_\theta(o_{i,t} | q, o_{i,<t})}{\pi_{\theta_{old}}(o_{i,t} | q, o_{i,<t})} \hat{A}_{i,t}, \,\, \text{clip} \left( \frac{\pi_\theta(o_{i,t} | q, o_{i,<t})}{\pi_{\theta_{old}}(o_{i,t} | q, o_{i,<t})}, 1 - \varepsilon, 1 + \varepsilon \right)  \hat{A}_{i,t} \bigg] - \beta \mathbb{D}_{KL}\left[\pi_{\theta} || \pi_{ref}\right]\bigg\} ,
\end{split}
\label{eq:GRPO-obj}
\end{equation}

where $\varepsilon$ and $\beta$ are hyper-parameters, and $\hat{A}_{i,t}$  is the advantage calculated based on relative rewards of the outputs inside each group only, which will be detailed in the following subsections. The group relative way that GRPO leverages to calculate the advantages, aligns well with the comparative nature of rewards models,  as reward models are typically trained on datasets of comparisons between outputs on the same question. Also note that, instead of adding KL penalty in the reward, GRPO regularizes by directly adding the KL divergence between the trained policy and the reference policy to the loss, avoiding complicating the calculation of  $\hat{A}_{i,t}$. And different from the KL penalty term used in (\ref{eq:PPO-reward}), we estimate the KL divergence with the following unbiased estimator \citep{kl_approx}:

\begin{equation}
\small
    \mathbb{D}_{KL}\left[\pi_{\theta} || \pi_{ref}\right] = \frac{\pi_{ref}(o_{i,t}|q,o_{i,<t})}{\pi_{\theta}(o_{i,t}|q,o_{i,<t})}- \log\frac{\pi_{ref}(o_{i,t}|q,o_{i,<t})}{\pi_{\theta}(o_{i,t}|q,o_{i,<t})} - 1,
\end{equation}

which is guaranteed to be positive.

\begin{algorithm}[t]
  \small
  \caption{Iterative Group Relative Policy Optimization}
  \textbf{Input} initial policy model $\pi_{\theta_{\text{init}}}$; reward models $r_\varphi$; task prompts $\mathcal{D}$; 
  hyperparameters $\varepsilon$, $\beta$, $\mu$
  \begin{algorithmic}[1]
    \State Set policy model $\pi_\theta$ to $\pi_{\theta_{\text{init}}}$
    \For{iteration = 1 to I}
       \State Assign reference model $\pi_{ref}$ as the current $\pi_{\theta}$
      \For{step = 1 to M}
      \State Draw batch $\mathcal{D}_b$ from the set $\mathcal{D}$
      \State Assign old policy $\pi_{\theta_{old}} \leftarrow \pi_{\theta}$
      \State For each question $q$ in $\mathcal{D}_b$, generate $G$ responses $\{o_i\}_{i=1}^G$ sampled from $\pi_{\theta_{old}} (\cdot \mid q)$
      \State Evaluate each generated output $o_i$ to compute rewards $\{r_i\}_{i=1}^{G}$ using $r_\varphi$
      \State For every output $o_i$, estimate token-wise group relative advantage $\hat{A}_{i,t}$
      \For{GRPO iteration = 1 to $\mu$}
        \State Optimize the policy $\pi_{\theta}$ by maximizing the Group Relative Policy Optimization (GRPO) objective (Equation \ref{eq:GC-GRPO})
      \EndFor
    \EndFor
    \State Continuously improve $r_\varphi$ utilizing replay-based training.
    \EndFor 
  \end{algorithmic}
  \textbf{Output} $\pi_\theta$
  \label{alg:iter-grpo}
\end{algorithm}

\subsubsection{Outcome Supervision RL with GRPO} In this approach, for each question $q$, a batch of outputs $\{o_1, o_2, \cdots, o_G\}$ is generated using the previous policy model $\pi_{\theta_{old}}$. A reward model assigns a corresponding set of rewards $\mathbf{r} = \{r_1, r_2, \cdots, r_G\}$ to these outputs. The obtained rewards are standardized by removing the mean of the group and dividing by their standard deviation. Under outcome supervision, the final normalized reward for each output $o_i$ is provided after completion, and this value is assigned as the advantage $\hat{A}_{i, t}$ to every token within $o_i$, i.e., $\hat{A}_{i, t} = \widetilde{r}_i = \frac{r_i- {\rm mean}(\mathbf{r})}{{\rm std}(\mathbf{r})}$. The policy is then updated by maximizing the objective specified in equation (\ref{eq:GRPO-obj}).

\subsubsection{Process Supervision RL with GRPO}  
While outcome supervision delivers a reward solely at the conclusion of each generated output, this sparse feedback can be insufficient for effectively guiding the policy, especially in intricate mathematical reasoning tasks. In line with the methodology described by \cite{wang2023math}, we extend our approach to include process supervision, wherein a reward is assigned at the completion of every intermediate reasoning step. Concretely, for a given question $q$ and $G$ sampled outputs $\{o_1, o_2, \cdots, o_G\}$, a process reward model evaluates each reasoning step in the outputs, generating rewards represented as $\mathbf{R} = \{ \{r_1^{index(1)},\cdots,r_1^{index(K_1)}\}, \cdots,  \{r_G^{index(1)},\cdots,r_G^{index(K_G)}\} \}$. Here, $index(j)$ denotes the token position corresponding to the termination of the $j$-th reasoning step, and $K_i$ signifies the number of reasoning steps for the $i$-th output. To standardize these reward values, we normalize them using their overall mean and standard deviation across all samples, giving $\widetilde{r}_i^{index(j)} = \frac{r_i^{index(j)} - {\rm mean(\mathbf{R})}}{{\rm std(\mathbf{R})}}$.

After normalization, the process supervision strategy computes the token-wise advantage by aggregating the normalized rewards from the current token onward, specifically: $\hat{A}_{i, t} = \sum_{index(j) \ge t} \widetilde{r}_i^{index(j)}$. The policy parameters are then updated to maximize the objective provided in equation (\ref{eq:GRPO-obj}).

\subsubsection{Iterative RL with GRPO} During the course of reinforcement learning, the existing reward model may become inadequate for effectively supervising the evolving policy model. To address this, we investigate iterative RL with GRPO. As illustrated in Algorithm \ref{alg:iter-grpo}, the iterative GRPO approach involves creating updated training datasets for the reward model, leveraging samples drawn from the current policy model. The reward model is further refined through ongoing training, utilizing a replay buffer that contains 10\% of past data. Subsequently, the reference model is updated to match the latest policy model, and training of the policy model resumes, now guided by the improved reward model.

\subsection{Training and Evaluating DeepSeekMath-RL}

We perform reinforcement learning utilizing the DeepSeekMath-Instruct 7B model. For RL training, we use approximately 144K chain-of-thought-style questions extracted from GSM8K and MATH within the SFT dataset. To assess how RL influences benchmarks where no relevant data is included during RL, we deliberately omit any other SFT questions. The reward model’s training dataset is prepared using the procedure described in \citep{wang2023math}. The initial reward model is trained using DeepSeekMath-Base 7B, setting the learning rate at 2e-5. For the policy model in the GRPO algorithm, we assign a learning rate of 1e-6 and use a KL penalty coefficient of 0.04. Each input prompt is used to sample $64$ candidate outputs. Both the maximum sequence length and the training batch size are set to 1024. The policy network undergoes just one update after each exploration cycle. To evaluate the resulting DeepSeekMath-RL 7B model, we use the same benchmarks as for DeepSeekMath-Instruct 7B. For DeepSeekMath-RL 7B, GSM8K and MATH tasks (with chain-of-thought solutions) are considered in-domain, whereas all remaining benchmarks are treated as out-of-domain.

Table \ref{tab:sft_rl_math} presents a comparative analysis of open-source and closed-source models employing both chain-of-thought and tool-augmented reasoning approaches on English and Chinese benchmark datasets. The results indicate the following: 1) DeepSeekMath-RL 7B achieves accuracy rates of 88.2\% on GSM8K and 51.7\% on MATH when using chain-of-thought reasoning, outperforming all other open-source models in the 7B to 70B parameter range and exceeding the accuracy of most closed-source counterparts. 2) Notably, DeepSeekMath-RL 7B's training involved only chain-of-thought-formatted instruction data from GSM8K and MATH, fine-tuned from the DeepSeekMath-Instruct 7B base. Although the range of training data was limited, it consistently surpasses DeepSeekMath-Instruct 7B on every evaluation metric, highlighting the substantial benefits conferred by reinforcement learning.

\section{Discussion}
This section presents the results and insights obtained from our investigations in both pre-training and RL experimental settings.

\subsection{Lessons Learnt in Pre-Training}

To begin, we describe our observations during the pre-training phase. Unless stated otherwise, the experimental configurations follow those detailed in Section \ref{sec:quality-policy}. Additionally, it should be emphasized that any mention of the DeepSeekMath Corpus in this section pertains to an 89B-token dataset generated in the second round of data collection.

\subsubsection{Code Training Benefits Mathematical Reasoning}

An often-cited but unproven assumption posits that learning to code enhances reasoning skills. In this work, we present a partial answer to this claim by focusing on mathematics: code training enhances models' performance in mathematical reasoning tasks, regardless of whether external tools are utilized.

To study how code training affects mathematical reasoning, we experimented with the following two-stage training and one-stage training settings:

\textbf{Two-Stage Training}

\begin{itemize}[topsep=0pt]
    \item \textbf{Code Training for 400B Tokens $\rightarrow$ Math Training for 150B Tokens}: DeepSeek-LLM 1.3B undergoes an initial training phase utilizing 400B code tokens, which is subsequently followed by a second stage focused on 150B math tokens.
    \item \textbf{General Training for 400B Tokens $\rightarrow$ Math Training for 150B Tokens}: To serve as a control, we conduct a parallel experiment where, during the first phase, the model is pretrained with 400B general tokens (extracted from DeepSeek-AI’s large-scale general-purpose corpus), then proceeds to the same 150B math token training. This setup enables us to compare the influence of code versus general token pretraining on the model’s mathematical reasoning performance.
\end{itemize}

\textbf{One-Stage Training}

\begin{itemize}[topsep=0pt]
    \item \textbf{Mathematics Pretraining for 150B Tokens}: DeepSeek-LLM 1.3B undergoes pretraining on 150B tokens sourced from math data;
    \item \textbf{Single-Stage Training Using 400B Code Tokens and 150B Math Tokens}: Sequential math training after code training diminishes the model's coding abilities. Thus, we examine the effects of jointly training on both code and math tokens in one unified stage, with the goal of not only enhancing mathematical reasoning but also mitigating catastrophic forgetting.
\end{itemize}

\paragraph{Results}
The downstream outcomes for various training configurations are presented in Table \ref{tab:code-math} and Table \ref{tab:code-math-general-eval}.

Program-aided mathematical reasoning is enhanced through code training, regardless of whether it is applied in a two-stage or one-stage training framework. As indicated in Table \ref{tab:code-math}, even employing only code training in the two-stage paradigm substantially improves performance on GSM8K and MATH tasks when leveraging Python. Incorporating a subsequent phase of math training delivers additional performance gains. Notably, within the one-stage training scheme, integrating both code and math tokens helps to prevent the catastrophic forgetting observed in the two-stage process, while simultaneously promoting improvements in both coding tasks (Table \ref{tab:code-math-general-eval}) and program-aided mathematical problem solving (Table \ref{tab:code-math}).

Code training enhances mathematical reasoning abilities even when tools are not employed. In a two-stage training framework, beginning with code training alone yields moderate improvements in mathematical performance. Moreover, this approach enhances the effectiveness of later math-specific training, culminating in superior results. In contrast, merging code and math tokens in a single training phase reduces the model's proficiency in mathematical reasoning absent tool support. A possible explanation is that DeepSeek-LLM 1.3B, given its relatively small size, does not possess the necessary capacity to comprehensively learn from both code and mathematics data at the same time.

\subsubsection{ArXiv Papers Seem Ineffective in Improving Mathematical Reasoning}

ArXiv papers are commonly included as a component of math pre-training data \citep{minerva,gpt-f,llemma,mathpile}. However, detailed analysis regarding their impact on mathematical reasoning has not been extensively conducted. Perhaps counter-intuitively, according to our experiments, arXiv papers seem ineffective in improving mathematical reasoning. We experiment with models of different sizes, including DeepSeek-LLM 1.3B and DeepSeek-Coder-Base-v1.5 7B \citep{deepseek-coder}, using arXiv corpora that underwent varied processing pipelines:

\begin{itemize}[topsep=0pt]
    \item \textbf{MathPile} \citep{mathpile}: a dataset consisting of 8.9 billion tokens, curated through a combination of heuristic filtering and data cleaning techniques, with scientific articles from arXiv comprising more than 85\% of the material;
    \item \textbf{ArXiv-RedPajama} \citep{redpajama}: this collection contains all arXiv LaTeX source files, with preambles, comments, custom macros, and bibliography sections omitted, amounting to 28.0 billion tokens in total.
\end{itemize}

In our study, we independently trained DeepSeek-LLM 1.3B on 150B tokens and DeepSeek-Coder-Base-v1.5 7B on 40B tokens, using each arXiv dataset in isolation. The results suggest that utilizing arXiv papers does not contribute to better mathematical reasoning abilities. Models exposed solely to arXiv-based data exhibit little to no progress—and, in some instances, performance even declines—across an array of mathematical benchmarks with varying levels of difficulty included in our evaluation. These assessments encompass quantitative reasoning datasets such as GSM8K and MATH (Table \ref{tab:arxiv-cot}), multiple-choice tests like MMLU-STEM (Table \ref{tab:arxiv-cot}), as well as formal mathematics benchmarks such as miniF2F (Table \ref{tab:arxiv_atp}).

However, this conclusion has its limitations and should be taken with a grain of salt. We have not yet studied:

\begin{itemize}[topsep=0pt]
    \item The influence of arXiv tokens on particular mathematics tasks outside the scope of this study, for example, transforming formal theorems or proofs into their informal counterparts (informalization);
    \item The interaction between arXiv tokens and additional data modalities;
    \item The extent to which the advantages conferred by arXiv papers persist or change as model size increases.
\end{itemize}

Thus, further exploration is required, which we leave for future studies.

\subsection{Insights of Reinforcement Learning}
\subsubsection{Towards to a Unified Paradigm} This section introduces a comprehensive framework for examining various training techniques, including SFT, RFT, DPO, PPO, and GRPO. We also carry out experiments to investigate the underlying elements of this unified framework. In general, the gradient with respect to the model parameter $\theta$ for any given training method can be formulated as:

\begin{equation}
    \nabla_{\theta}\mathcal{J}_{\textcolor{red}{\mathcal{A}}}(\theta) = \mathbb{E}[\underbrace{(q,o) \sim \textcolor{red}{\mathcal{D}}}_{Data \ Source}]\left( \frac{1}{|o|} \sum_{t=1}^{|o|}  \underbrace{GC_{{\mathcal{A}}}(q, o, t, \textcolor{red}{\pi_{{rf}}})}_{Gradient \ Coefficient}  \nabla_{\theta}\log \pi_{\theta}(o_t | q, o_{<t})\right).
\label{eq:objective}
\end{equation}

There exist three key components: 1) \textit{Data Source $\mathcal{D}$}, which determines the training data; 2) \textit{Reward Function $\pi_{{rf}}$}, which is the source of the training reward signal; 3) \textit{Algorithm $\mathcal{A}$}: which processes the training data and the reward signal to the gradient coefficient $GC$ that determines the magnitude of the penalty or reinforcement for the data. We analyze several representative methods based on such a unified paradigm:

\begin{itemize}
    \item  \textbf{Supervised Fine-tuning (SFT)}: In SFT, a pretrained model is further adapted by training it on a dataset curated by humans.
    \item \textbf{Rejection Sampling Fine-tuning (RFT)}: RFT enhances the SFT-trained model by additionally training it on those responses sampled from the SFT model which pass a correctness-based filtering criterion tied to the SFT dataset.
    \item \textbf{Direct Preference Optimization (DPO)}: DPO builds upon the SFT model by training with pairs of augmented responses drawn from the SFT model, applying a pairwise DPO loss to improve performance.
    \item \textbf{Online Rejection Sampling Fine-tuning (Online RFT)}: Unlike standard RFT, Online RFT starts with the SFT model as the policy and updates it through fine-tuning with augmented samples obtained from the current policy model in an online manner.
    \item \textbf{PPO/GRPO}: PPO/GRPO initializes with the SFT-trained policy model and then strengthens the policy through reinforcement learning, leveraging outputs sampled from the up-to-date policy during training.
\end{itemize}

A summary of the elements comprising these methods is provided in Table \ref{tab:data_policy}. For a comprehensive explanation of the derivation steps, see Appendix \ref{app:analysis-rl}.

\paragraph{Observation about Data Source} We categorize the data source into two types: online sampling and offline sampling. Online sampling refers to gathering training data from the outcomes produced by the actively learning policy model during training, whereas offline sampling utilizes training data derived from the sampling outputs of the pretrained SFT model. Both RFT and DPO adopt the offline sampling approach, whereas Online RFT and GRPO rely on the online sampling method.

As illustrated in Figure \ref{fig:rl-analysis}, our results indicate that Online RFT consistently surpasses RFT across both benchmarks. Notably, the performance of Online RFT mirrors that of RFT during the early phases of training; however, as training advances, Online RFT establishes a clear lead, highlighting the benefits of online training methods. This observation aligns with expectations: at the outset, both the actor and the SFT model behave similarly, so the collected samples show minimal variation. Conversely, as training proceeds, the divergence between the data generated by the actor and that from the SFT model becomes more pronounced, at which point the advantages of sampling data in real time become more apparent.

\paragraph{Observation about Gradient Coefficient} In updating the model parameters, the algorithm employs the gradient coefficient derived from the processed reward signal. In our experimental setup, we distinguish the reward function into two components: `Rule' and `Model'. The `Rule' component evaluates the response quality strictly according to answer correctness, whereas the `Model' component involves training a reward model that assigns a score to each response. The training set for the reward model is constructed based on the initial rule-based evaluations. As presented in Equations \ref{eq:GC-RFT} and \ref{eq:GC-GRPO}, there is an important distinction between the GRPO and Online RFT approaches: GRPO uniquely tailors the gradient coefficient according to the reward model’s output for each response, thereby enabling nuanced reinforcement or penalization aligned with the specific reward magnitude. Conversely, Online RFT does not incorporate such granularity; it applies uniform reinforcement to all correctly answered responses and refrains from penalizing incorrect ones.

Figure \ref{fig:rl-analysis} illustrates that GRPO outperforms online RFT, underscoring the effectiveness of modifying the coefficients associated with positive and negative gradients. Moreover, the results reveal that GRPO+PS achieves better outcomes than GRPO+OS, emphasizing the advantage of incorporating more precise, step-sensitive gradient coefficients. Additionally, our investigation extends to iterative RL procedures, where we carry out two iterations during experimentation. As presented in Figure \ref{fig:iter-rl}, the implementation of iterative RL yields notable performance gains, most prominently during the initial iteration.

\subsubsection{Why RL Works?} In this study, we apply reinforcement learning using a subset of instruction tuning datasets, resulting in notable gains over the baseline instruction tuning approach. To elucidate the mechanism underlying the effectiveness of reinforcement learning, we compare the Pass@K and Maj@K accuracies for both the Instruct and RL models across two evaluation datasets. Figure \ref{fig:maj-pass} illustrates that reinforcement learning yields better Maj@K performance, while Pass@K remains largely unchanged. This suggests that reinforcement learning primarily improves the model by increasing the stability of its output distribution; in effect, \textbf{the observed gains appear to stem from elevating the most likely correct response within TopK candidates, rather than advancing the model’s core skill set.} In a related vein, \citep{wang2023making} discuss a \textbf{misalignment problem} present in reasoning tasks for SFT models and demonstrate that applying preference alignment techniques \citep{yuan2023rrhf,song2023preference,wang2023making} can successfully enhance reasoning performance.

\subsubsection{How to Achieve More Effective RL?}
\label{sec:effec_rl}
Our experiments indicate that RL is highly effective for mathematical reasoning tasks. Additionally, we introduce a cohesive framework that clarifies the relationships among various prominent training strategies. In this framework, each approach can be interpreted as either an explicit application or an approximation of RL methods. As articulated in Equation \ref{eq:objective}, there are three crucial elements: Data Source, Algorithm, and Reward Function. We outline possible avenues for advancing each of these components in future research.

\paragraph{Data Source} The data source forms the foundation for all training methodologies. Within the framework of RL, we define the data source as unlabeled questions alongside the responses generated by the policy model. For this work, we restrict ourselves to using questions originating from the instruction tuning phase, and we generate responses by applying a straightforward nucleus sampling strategy. This limited approach may explain why our RL process primarily enhances Maj@K performance metrics. Going forward, we intend to assess our RL pipeline’s effectiveness on prompts that lie outside the distribution of the training data, pairing this with \textbf{advanced sampling (decoding) strategies} such as those leveraging tree-search methodologies \citep{yao2023tree}. Furthermore, \textbf{efficient inference techniques} \citep{xia-etal-2023-speculative,leviathan2023fast,kwon2023efficient,xia2024unlocking}, which are critical to improving exploration within policy models, are expected to have a significant impact as well.

\paragraph{Algorithms} Algorithms utilize data and the reward signal to compute the gradient coefficients needed for model parameter updates. According to Equation \ref{eq:objective}, existing approaches generally place full \textbf{TRUST} in the feedback provided by the reward function, leveraging it to either increase or decrease the conditional probability assigned to particular tokens. Nevertheless, the reliability of reward signals cannot be universally guaranteed, particularly in highly complex scenarios. For instance, the PRM800K datasets \citep{lightman2023let}, despite being meticulously annotated by expert annotators, still present a misannotation rate of approximately 20\%\footnote{\url{https://github.com/openai/prm800k/issues/12\#issuecomment-1728491852}}. Therefore, we intend to investigate reinforcement learning methods designed to withstand noisy or inaccurate reward signals. We argue that alignment strategies characterized as \textbf{WEAK-TO-STRONG} \citep{burns2023weak} could substantially transform current learning paradigms.

\paragraph{Reward Function} The reward function serves as the primary source of training feedback. Within RL, this function typically takes the form of a neural reward model. We identify three crucial avenues for the development of reward models: 1) \textbf{Improving the generalization capacity of the reward model.} It is essential for the reward model to generalize effectively to novel, out-of-distribution prompts and more sophisticated decoding outputs; in the absence of such generalization, reinforcement learning may only serve to consolidate the existing distribution of LLMs, failing to advance their underlying abilities; 2) \textbf{Capturing and leveraging the uncertainty inherent in reward models.} Modeling uncertainty may offer a way to connect unreliable reward models with algorithms designed for weak-to-strong learning; 3) \textbf{Constructing high-quality, process-level reward models with efficiency,} such that they can deliver granular feedback during multi-step reasoning tasks \citep{lightman2023let,wang2023math}.

\section{Conclusion, Limitation, and Future Work}

In this work, we introduce DeepSeekMath, which achieves superior performance compared to all other open-source models on the MATH benchmark at the competition level and nearly matches the effectiveness of proprietary solutions. DeepSeekMath~builds upon DeepSeek-Coder-v1.5 7B as its foundation, and is further trained continuously on 500B tokens, incorporating 120B math-related tokens mainly extracted from Common Crawl. Through an in-depth ablation analysis, we find that content from web sources represents a rich reservoir for high-quality mathematical datasets, whereas arXiv does not contribute as much as anticipated. We propose Group Relative Policy Optimization (GRPO), a new adaptation of Proximal Policy Optimization (PPO), which enhances mathematical reasoning with improved memory efficiency. Experimental evidence indicates that, even when DeepSeekMath-Instruct 7B already attains high benchmark results, GRPO leads to additional gains. Furthermore, we outline a comprehensive framework for interpreting a family of reinforcement learning strategies and highlight several promising avenues for developing more effective techniques in this area.

While DeepSeekMath~demonstrates strong performance on benchmarks assessing quantitative reasoning, its abilities in geometry and theorem-proving tasks fall short compared to proprietary models. For example, in our preliminary evaluation, the model was unable to solve questions involving triangles and ellipses, which could point to pre-training and fine-tuning data selection biases. Furthermore, due to the model’s limited size, DeepSeekMath~exhibits inferior few-shot learning capacity relative to GPT-4. Whereas GPT-4 benefits from few-shot prompts, DeepSeekMath~displays no significant difference between zero-shot and few-shot scenarios. Moving forward, we plan to enhance our data selection workflow to develop a more robust, high-quality pre-training dataset. Additionally, we aim to investigate the possible avenues for improved reinforcement learning strategies for LLMs, as outlined in Section \ref{sec:effec_rl}.

\end{document}